% !TeX root = CV_Tramarin.tex

%\section{Teaching support activities and Professional Experience}
%% COLLABORAZIONI

%\section{Research Collaborations and Consulting}

    %\textbf{\textit{Trasferimento tecnologico e contratti di ricerca}}
    %
    %Le attività di ricerca descritte, coerentemente anche con le competenze
    %acquisite sui sistemi embedded per applicazioni industriali, hanno avuto
    %rilevanza nell’ambito delle attività di trasferimento tecnologico e di contratti
    %di ricerca con alcune realtà industriali del Triveneto. Un esempio significativo
    %riguarda la collaborazione con MC Elettronica Srl, azienda che progetta e
    %sviluppa sistemi embedded per l’agricoltura di precisione. Nel dettaglio, la
    %collaborazione ha riguardato lo sviluppo di un sistema embedded da installare a
    %bordo macchina, che permettesse l'acquisizione remota di dati di misura dal
    %macchinario, permettendone un monitoraggio  real--time remoto dello stato.
    %Inoltre, il sistema oggetto di studio doveva anche permettere un’eventuale
    %attuazione remota di sezioni non critiche tramite una applicazione specifica da
    %eseguire in un comune dispositivo (tablet, PC o smartphone) commerciale
    %\cite{etfa14_a}.
    %
    %Sono stato inoltre co-responsabile di un contratto di ricerca con Enerspin Srl,
    %azienda che produce distributori rotanti per macchinari industriali. L'attività
    %connessa a questa collaborazione è legata alla progettazione e allo sviluppo di
    %un banco di prova per la validazione delle linee di trasferimento dati che
    %attraversano il distributore rotante. In particolare questo banco di prova
    %doveva permettere l'esecuzione di test di pre--compliance dei distributori
    %prodotti rispetto alle performance attese da specifici protocolli di
    %comunicazione industriali real--time Ethernet \cite{imtc2018}.
    %
    %Inoltre, in ambito di trasferimento tecnologico e di saperi, sono stato nominato
    %membro del Comitato Scientifico della Fiera SPS (Smart Production Systems), la
    %più rilevante in ambito italiano e europeo (gemellata con l'omonima fiera
    %annuale di Norimberga) per il settore dell'automazione industriale. In questo
    %contesto sono stato invitato, come esperto italiano, a tenere un intervento alla
    %Tavola Rotonda inaugurale dell'edizione 2019 per delineare lo stato di
    %avanzamento della standardizzazione dei sistemi Time Sensitive Networking (TSN),
    %che rappresenteranno il futuro delle comunicazioni in ambito industriale e dei
    %sistemi di misura distribuiti per il controllo di processo. Sulle stesse
    %tematiche ho anche tenuto un intervento su invito in una recente Tavola Rotonda,
    %a cui hanno partecipato molti attori dello scenario Industria 4.0, come Siemens,
    %Rockwell, Hitachi per la parte di automazione industriale, e Qualcomm, Tim e
    %Vodafone per quanto riguardo lo sviluppo del 5G in queste applicazioni.

\sectioniten{Responsabilità di studi e ricerche affidate da istituzioni pubbliche o private qualificate}{Responsibility for scientific studies and research assigned by qualified public or private institutions}

% Utilizzo della tecnologia
% RFID per il tracciamento degli assistiti in strutture sanitarie”. L’attività si colloca nell'ambito del
% progetto di ricerca finanziato dalla Regione Veneto, nell’ambito del Bando per il sostegno a progetti di
% ricerca e sviluppo realizzati da Aggregazioni di imprese sul tema operatore sanitario supportato da
% tecnologie digitali denominato “Hospice Operator Digital Twin – HODT”

\cvitemit{2023-2024}{{\sbseries Responsabilità scientifica del contratto di Ricerca tra CNR-IEIIT e Dataveneta Srl.}\newline
    Utilizzo della tecnologia RFID per il tracciamento degli assistiti in strutture sanitarie.\newline L’attività si colloca nell'ambito del progetto di ricerca finanziato dalla Regione Veneto, realizzati da Aggregazioni di imprese sul tema operatore sanitario supportato da tecnologie digitali denominato “Hospice Operator Digital Twin (HODT)}

\cvitemen{2023-2024}{{\sbseries Research Contract between CNR-IEIIT and Dataveneta Srl.}\newline
    Analysis of innovative systems for interoperability and advanced connectivity to be applied to embedded systems for automation and distributed measurements in industry based on OPC UA}

\cvitemit{2018}{Contratto di Ricerca tra Università di Padova e CMZ Srl.\newline
    Analisi di sistemi innovativi per l'interoperabilità e la connettività avanzata da applicare a sistemi embedded per l'automazione e le misure distribuite nell'industria basati su OPC UA.}

\cvitemen{2018}{Research Contract between University of Padova and CMZ Srl.\newline
    Analysis of innovative systems for interoperability and advanced connectivity to be applied to embedded systems for automation and distributed measurements in industry based on OPC UA}

\cvitemit{2017}{{\sbseries Co-Responsabile Scientifico, Contratto di Ricerca, tra CNR--IEIIT ed Enerspin Srl.}\newline
    Sviluppo e implementazione di un sistema embedded per la valutazione e la validazione dell'integrità dei dati nei distributori rotanti per applicazioni industriali.}

\cvitemen{2017}{{\sbseries Scientific Co-Responsible, Research Contract, between CNR--IEIIT and Enerspin Srl.}\newline
    February 2017\newline
    Development and implementation of an embedded system for the assessment and validation of data integrity in rotary distributors for industrial applications.}

\cvitemit{2016}{Contratto di Ricerca tra CNR--IEIIT e KBlue Srl\newline
Protocollo n. IEIIT 1090 --- 27/07/2016\newline
Nell'ambito dei sistemi domotici, analisi delle prestazioni di sistemi distribuiti basati su interconnessioni powerline.}

\cvitemen{2016}{Research Contract between the CNR--IEIIT and KBlue Srl\newline
Protocol no. IEIIT 1090 --- 27/07/2016\newline
Within the field of domotic systems, performance analysis of distributed systems based on powerline interconnections.}

\cvitemit{2016}{Contratto di Ricerca tra CNR--IEIIT e MC Elettronica Srl\newline
Protocollo n. IEIIT 1091 --- 27/07/2016\newline
Ricerca e implementazione di un sistema embedded per il controllo remoto e la supervisione di macchinari agricoli con sistemi wireless basati su IEEE 802.11g.}

\cvitemen{2016}{Research Contract between the CNR--IEIIT and MC Elettronica Srl\newline
Protocol no. IEIIT 1091 --- 27/07/2016\newline
Research and implementation of an embedded system for remote control and supervision of machinery in agricolture with wireless IEEE 802.11g based systems.}

%\cvitem{Feb--Mar 2014}{Research Collaboration with the Electronic
%Measurement Research Group, Dept. Information Engineering, University of Padova, for the project ``Modeling, simulations and experimental assessment of the performance of industrial protocols in smart sensor networks''.}
%
%\cvitem{June 2012}{Research and technical collaboration, Research Contract, Electronic
%Measurement Research Group (University) and Silverstar Engineering
%Srl, consulting on the analysis and suitable techniques for the positioning and
%effective setup of wireless sensors, for building structural monitoring.}

%\cvitem{Apr 2009--June 2010}{Research and technical collaboration, Research Contract, Electronic
%Measurement Research Group (University of Padova) and Carel
%Industries SpA, consulting on high-speed ARM-based electronic design, in particular
%to solve EMC related issues (signal integrity, radiated and conducted emissions).}
%
%\cvitem{Oct 2007--Dec 2008}{Research collaboration with the Electronic 
%Measurement Research Group (University of Padova), started during the master
%degree thesis activities and till the start of the PhD}
%
%\section{Thesis supervision}
%
%\cvitem{Supervisione di dottorandi}{Supervisione di 1 progetto di Tesi Dottorale per il Corso di dottorato di Ricerca in Ingegneria Meccatronica e dell’Innovazione Meccanica del Prodotto, Dipartimento di Tecnica e Gestione dei Sistemi Industriali (DTG), Università degli Studi di Padova}
%
%\cvitem{Supervisione di tesi}{Relatore di più di 10 Tesi di Laurea Triennale e Magistrali}
%
%\cvitem{Controrelazioni di tesi}{Controrelatore di 3 Tesi di Laurea Magistrali}
%
%\cvitem{Co-Supervisione di tesi}{Correlatore oltre 20 Tesi di Laurea Triennale e Magistrali}

% \cvitem{Supervisor of Bachelor Thesis}{``Analisi delle prestazioni di un link wireless IEEE 802.11 per l’interconnessione tra due segmenti real-time ethernet industriali'', E. Petri, Bachelor degree in  Mechanical and Mechatronic Enginering.  
% Thesis defense: 28th november 2017.}

% \cvitem{Supervisor of Bachelor Thesis}{``Programmazione di Dispositivi Industriali per la misura delle prestazioni di reti real-time Profibus DP'', R. Salviati, Bachelor degree in  Mechanical and Mechatronic Enginering.  
% Thesis defense: 28th november 2017.}

% \cvitem{Supervisor of Bachelor Thesis}{``Programmazione di Dispositivi Industriali per il Collegamento a Reti Real-Time Profibus DP'', M. Simioni, Bachelor degree in  Mechanical and Mechatronic Enginering.  
% Thesis defense: 18th september 2017.}

%\cvitem{Co-supervisor of Bachelor thesis}{``Stima di Canale per Autenticazione mediante 
%Software Defined Radio (Channel Estimation through Software Defined Radio for Authentication)'', 
%A. Cavicchio, Bachelor degree in Computer Science. 
%Thesis period: 9 months. 18 March 2015.}
%
%\cvitem{Co-supervisor of Bachelor thesis}{``Powerline Communication per sistemi
%Domotici (Powerline communication for Home Automation Systems)'', 
%M. Dominese, Bachelor degree in Automation Engineering. 
%Thesis period: 3 months. Data discussione: 18 March 2015.}
%
%\cvitem{Co-supervisor of Master thesis}{``Traffic Models for Data Networks and Energy 
%Efficient Control'', 
%M. Luvisotto, Master degree in Automation Engineering. 
%Durata tesi: 6 mesi. 7 July 2014.}
%
%\cvitem{Co-supervisor of Bachelor thesis}{``Implementazione dello standard IEEE 802.11g su sistemi 
%Software Defined Radio (IEEE 802.11g implementation on Software Defined Radio)'', 
%F. Casamento, Bachelor degree in Information Engineering. 
%Thesis period: 9 months. 20 February 2014}
%
%\cvitem{Co-supervisor of Master thesis}{``Traffic Models for Data Networks and Energy 
%Efficient Control'', 
%M. Michielan, Master degree in Automation Engineering. 
%Thesis period: 6 months. 9 December 2013}
%
%\cvitem{Co-supervisor of Bachelor thesis}{``Reti ethernet a efficienza energetica per applicazioni industriali. Parte A: lo standard 802.3az (Energy Efficient Ethernet networks for industrial application. Part A: the IEEE 802.3az standard)'', 
%M. Conti, Bachelor degree in Information Engineering.
%Thesis period: 3 months. 30 September 2013}
%
%\cvitem{Co-supervisor of Bachelor thesis}{``Reti ethernet a efficienza energetica per applicazioni industriali. Parte B: le misure in laboratorio (Energy Efficient Ethernet networks for industrial application. Part B: experimental measurements)'', 
%M. Perazzolo, Bachelor degree in Information Engineering. 
%Thesis period: 3 months. 30 September 2013}
%
%\cvitem{Co-supervisor of Master thesis}{``Ottimizzazione del consumo energetico per reti
%Ethernet Real-Time (Optimization of energy consumption for ethernet networks)'', 
%A. Salata, Master degree in Automation Engineering.
%Thesis period: 6 months. March 11th, 2013}
%
%\cvitem{Co-supervisor of Master thesis}{``Studio e progetto di metodi di misura robusti
%per sincrofasori in bassa tensione (Design of robust measurement methods for synchrophasors in low voltage lines)'', 
%M. Penzo, Master degree in Electronic Engineering.
%Thesis period: 9 months. October, 2012}
%
%\cvitem{Co-supervisor of Bachelor thesis}{`Impiego di un dispositivo basato su FPGA per misure
%di power quality in Real-Time (Real--time power quality measurements through an FPGA device)'', 
%I. Ergest, Bachelor degree in Electronic Engineering.
%Thesis period: 6 months. February 25th, 2013}
%
%\cvitem{Co-supervisor of Bachelor thesis}{``Simulazione cross-layer di tecniche di ritrasmissione
%in reti wireless interferenti 802.11 e 802.15.4 (Cross-layer simulation of retransmission strategies for interfering IEEE 802.11 and IEEE 802.15.4 networks)'', 
%R. Bersan, Bachelor degree in Electronic Engineering.
%Thesis period: 6 months. February 21st, 2011}
%

%\subsection{Attività di Docenza in Corsi di Formazione e post Diploma}
%
%\cvitem{2009-- In corso}{Regolarmente impegnato come Docente in corsi ITS, soprattutto nel comparto Meccatronico}
%
%\cvitem{2009-- In corso}{Attività formative di carattere aziendale --- Docente in Corsi di formazione e approfondimento all'interno di collaborazioni Università--Imprese, soprattutto per gli aspetti di\newline
%\parbox{\textwidth}{
%\begin{itemize}
%\item Sensoristica per applicazioni Industriali
%\item Reti Ethernet Industriali (RTE)
%\item Compatibilità elettromagnetica
%\end{itemize}
%}}

%% INSEGNAMENTO

%\subsection{Teaching activities}
%
%\cvitem{Feb 2014 -- June 2015}{Lecturer and Tutor for Team Working activities for the development of prototypal automation device, within the ITS Meccatronico School (High Technology Institute for Mechatronics), Padova, Italy, 80 hours.}
%
%\cvitem{Dec 2014--Jan 2015}{Lecturer for the Course “Reti per l'Automazione” (Communication Networks for Industrial Automation) within the ITS Meccatronico School (High Technology Institute for Mechatronics), Padova, Italy, 12 hours.}
%
%\cvitem{Nov--Dec 2014}{Lecturer for the Course “Fundamentals of Computer Science” within the ITS Cosmo School (High Technology Institute for Design), Padova, Italy, 20 hours.}
%
%\cvitem{Dec 2013--Jan 2014}{Lecturer for the Course “Reti per l'Automazione” (Communication Networks for Industrial Automation) within the ITS Meccatronico School (High Technology Institute for Mechatronics), Padova, Italy, 20 hours.}
%
%\cvitem{II semester, A.Y. 2012/2013}{Teaching assistant, laboratory activities,
%Course of Industrial Automation, Master Degree in Automation Engineering.}
%
%\cvitem{I semester, A.Y. 2012/2013}{Lectures, Electromagnetic Compatibility
%Course, Master Degree in Electronic Engineering, 6 hours.\newline 
%Lectures, Electronic Instrumentation
%Course, Degree in Biomedic Engineering, 4 hours.}
%
%\cvitem{II semester, A.Y. 2011/2012}{Teaching assistant and workgroup tutor,
%Course on Smart Grid, Master Degree in Information Engineering.}
%
%\cvitem{I semester, A.Y. 2011/2012}{Lectures, Electromagnetic Compatibility
%Course, Master Degree in Electronic Engineering, 6 hours.}
%
%\cvitem{Feb--May 2011}{Lecturer for the Course “Tecnico per la certificazione
%nell'industria elettronica” (Technician for the Certification in the Electronic 
%Industry), Technical Institute ITIS F. Severi, Padova, Italy, 30 hours.}
%
%\cvitem{March--April 2011}{Lecturer for the Course “Metodologie di misurazione, norme e
%progettazione di apparecchiature elettroniche finalizzate alla marcatura CE"
%(Methodologies for measurement, design and regulations related to
%electronic systems, for the CE compliance) -- in
%collaboration with TIS Innovation Park, and the Autonomous Province of Bolzano,
%Italy, 24 hours.}
%
%\cvitem{March--May 2010}{Lecturer for the Course “Metodi e strumenti di misura per prove
%di compatibilità elettromagnetica finalizzate alla marcatura CE” 
%(Methods and instrumentation for electromagnetic compatibility testing
%for the CE compliance)-- in
%collaboration with TIS Innovation Park, and the Autonomous Province of Bolzano,
%Italy, 32 hours.}
%
%\cvitem{Apr--June 2010\newline Sept--Oct 2009}{Lecturers, training Course for Electronic Designers (engineers), Carel Industries S.p.A, 8 hours.}
%
%\cvitem{A.Y. 2009/2012}{Teaching assistant, laboratory activities, Courses of
%Electronic Measurement, Electromagnetic Compatibility, Electronic Instrumentation,
%collaboration contract with the Engineering Faculty of the university of
%Padova, Italy.}
%
%\cvitem{A.Y. 2007/2010}{Tutor for students of the Faculty of Engineering, Faculty of Engineering,University of
%Padova,\newline - teaching assistant for laboratory activities of the Courses of
%Fundamentals of Programming (I year, I semester); \newline - lecturers for students
%of the first year Course in Mathematical Analysis 1 (guided exercises); \newline - lectures and organization of a course on Fundamentals of \LaTeX;
%\newline - teaching assistant for laboratory activities, Courses of Electronic Measurement, Electromagnetic Compatibility, Electronic Instrumentation.}



